\documentclass[11pt]{letter}

\usepackage[margin=0.5in]{geometry}
\usepackage{hyperref}

\begin{document}

\begin{letter}{48Hour Discovery}

\opening{Dear Hiring Team,}

I am writing to apply for the Data Development Intern position at 48Hour Discovery. I recently completed my M.Sc. in Computer Science at the University of Calgary, and I am interested in this role because it combines scientific data, Python development, web services, databases, cloud/Linux workflows, and practical support for research teams. I am based in Alberta and would be excited to contribute to data infrastructure that helps accelerate discovery work.

My graduate research focused on building Python and C++ software for large scientific and geospatial datasets. I developed pipelines for data collection, preprocessing, validation, storage, retrieval, analysis, model training, and performance evaluation. One major project involved a DGGS-based framework for managing large spatiotemporal satellite datasets, where I implemented multiresolution storage and retrieval methods and evaluated the system on large Canada-wide data. This gave me strong practice working with complex datasets, designing reliable data flows, and documenting technical decisions clearly.

I also worked on a digital elevation model super-resolution project using Python, NumPy, pandas, Rasterio, Google Earth Engine, and PyTorch. I built workflows to collect and preprocess terrain data, trained deep learning models, and evaluated outputs with statistical and domain-specific metrics such as RMSE, slope, TRI, TPI, and wavelet-based errors. That project strengthened my ability to debug data pipelines, validate results, and think carefully about how data quality affects downstream analysis.

In addition to research software, I bring backend experience from Asr Gooyesh Pardaz, where I developed Django REST APIs backed by PostgreSQL for NLP, text-to-speech, and speech-to-text products. I worked on authentication features, payment-service integration, relational records, and API behavior used by production services. While the 48Hour Discovery role mentions Flask specifically, my Django and REST API experience gives me a strong foundation for contributing to Python web tooling and learning the team's framework conventions quickly.

I am also comfortable working in Linux environments, using Git, writing documentation, and communicating technical ideas to different audiences. As a teaching assistant in database systems, data-at-scale, big-data applications, and deep learning, I helped students debug code, reason about data workflows, and move projects from unclear requirements to working implementations. I have less direct biotechnology and R experience, but I bring strong Python, SQL, scientific-data, cloud-data, and ML validation experience, along with the curiosity and discipline needed to learn the domain carefully.

Thank you for considering my application. I would welcome the opportunity to discuss how my Python, data-pipeline, backend, and research-software background could support 48Hour Discovery's data development work.

\closing{Sincerely,

Amir Mirzai Golpayegani

\href{mailto:amirgolpaa24@gmail.com}{amirgolpaa24@gmail.com}}

\end{letter}

\end{document}
